\documentclass[12pt,a4paper]{article}
\usepackage{mathwizards-ingles}
\title{%
  \vspace{-1cm}
  {\Huge\bfseries Cronograma de Sesiones Sincrónicas}\\[0.3cm]
  {\Large 36 Encuentros Virtuales · 12 Semanas · 3 Meses}\\[0.2cm]
  {\large\color{grissuave} \textcolor{mes1}{$\blacksquare$ Mes 1} \quad \textcolor{mes2}{$\blacksquare$ Mes 2} \quad \textcolor{mes3}{$\blacksquare$ Mes 3}}\\[0.5cm]
  \rule{\textwidth}{1.5pt}
}
\author{}
\date{}

\begin{document}
\maketitle
\thispagestyle{empty}

% ═══════════════════════════════════════════════════════════════════════════════
% MES 1
% ═══════════════════════════════════════════════════════════════════════════════

\begin{cajames}[\textcolor{white}{Mes 1: El Periodo Silencioso y Adquisición Fonológica (A0 $\rightarrow$ A1)}]{mes1}
\textbf{Dinámica:} El profesor actúa, gesticula y usa material visual. Los estudiantes responden sin forzar el habla (reacciones no verbales, emojis, TPR).
\end{cajames}

\vspace{0.3cm}

\begin{cajasemana}[\textcolor{white}{Semana 1: Adaptación y Rompehielo}]{mes1}
\sesion{1}{Bienvenida y Setup Técnico}{Configuración de Anki, entrega del mazo base grupal. Dinámica rompehielo de ``señalar y reaccionar'' (cero inglés forzado).}
\sesion{2}{TPR Básico (Respuesta Física Total)}{El profesor modela verbos de acción y objetos de la habitación; los alumnos reaccionan en cámara.}
\sesion{3}{Watch-Party Guiada}{Visualización grupal de video \textit{Absolute Beginner}. Pausas para asegurar que se infiere el contexto global.}
\end{cajasemana}

\begin{cajasemana}[\textcolor{white}{Semana 2: Expansión Visual}]{mes1}
\sesion{4}{TPR Interactivo}{Cuerpo, emociones, direcciones. Uso de pizarras virtuales (whiteboards).}
\sesion{5}{Inmersión Narrativa ``Slice of Life''}{Análisis grupal de un videoblog (rutina matutina) con pausas estratégicas.}
\sesion{6}{Dinámica de Comprensión}{El profesor narra una mini-historia con soporte visual. Alumnos responden encuestas rápidas (polls).}
\end{cajasemana}

\begin{cajasemana}[\textcolor{white}{Semana 3: Consolidación del Contexto}]{mes1}
\sesion{7}{Cuentacuentos (ALG)}{Historia sencilla apoyada 100\% en dibujos en vivo (Crecimiento Automático del Lenguaje).}
\sesion{8}{Role-play Mudo}{El profesor actúa escenarios cotidianos; los estudiantes eligen el curso de la historia votando en el chat.}
\sesion{9}{Watch-Party}{Video de \textit{Comprehensible English}. Discusión en foro (en español) sobre las emociones o el mensaje del video.}
\end{cajasemana}

\begin{cajasemana}[\textcolor{white}{Semana 4: Transición y Herramientas}]{mes1}
\sesion{10}{Repaso Lúdico Visual}{Kahoot/Quizizz solo con imágenes y audio nativo, sin texto.}
\sesion{11}{Historia A1 Narrada en Vivo}{Integración del vocabulario de alta frecuencia del mazo Anki del grupo.}
\sesion{12}{Taller Técnico Preparatorio}{Instalación en vivo de extensiones (Language Reactor/Yomitan) y AnkiConnect. Capacitación sobre la regla ``1T'' (\textit{i+1}) para el Mes 2.}
\end{cajasemana}

% ═══════════════════════════════════════════════════════════════════════════════
% MES 2
% ═══════════════════════════════════════════════════════════════════════════════

\newpage

\begin{cajames}[\textcolor{white}{Mes 2: Minería Colaborativa y Lectura Asistida (A1 $\rightarrow$ A2)}]{mes2}
\textbf{Dinámica:} Lectura compartida e interacción analítica con el vocabulario. El profesor guía la inferencia colaborativa.
\end{cajames}

\vspace{0.3cm}

\begin{cajasemana}[\textcolor{white}{Semana 5: Introducción a la Lectura}]{mes2}
\sesion{13}{Círculo de Lectura Virtual I}{El profesor proyecta un \textit{Graded Reader} (Nivel A1). Lectura compartida con audio nativo.}
\sesion{14}{Círculo de Lectura Virtual II}{Continuación del texto. Ejercicios de inferencia de contexto antes de revelar el significado.}
\sesion{15}{Taller Práctico de Minería}{Compartir pantalla; los alumnos presentan sus primeras oraciones minadas. El profesor valida \textit{i+1}.}
\end{cajasemana}

\begin{cajasemana}[\textcolor{white}{Semana 6: Independencia Contextual}]{mes2}
\sesion{16}{Círculo de Lectura III}{Enfoque en tiempos verbales en contexto narrativo (pasados/futuros).}
\sesion{17}{Análisis Narrativo Visual}{Consumo de un video nivel A2 (\textit{EnglishSponge}).}
\sesion{18}{Sesión de \textit{Crowdsourcing} de Tarjetas}{Los estudiantes comparten sus mejores oraciones 1T en Discord; el profesor aprueba y cura un mini-mazo semanal grupal.}
\end{cajasemana}

\begin{cajasemana}[\textcolor{white}{Semana 7: Lectura Extensiva Guiada}]{mes2}
\sesion{19}{Círculo de Lectura IV}{Cambio de libro o género (misterio/aventura nivel A1/A2).}
\sesion{20}{Lectura a Vista}{El profesor muestra párrafos no preparados. Lectura ágil fomentando tolerar palabras desconocidas.}
\sesion{21}{Clínica de Anki}{Revisión de malos hábitos (tarjetas defectuosas, traducciones en el lado frontal) y erradicación de ``sanguijuelas''.}
\end{cajasemana}

\begin{cajasemana}[\textcolor{white}{Semana 8: Transición Semántica}]{mes2}
\sesion{22}{Círculo de Lectura Final del Mes}{Conclusión del libro del mes. Resumen colaborativo.}
\sesion{23}{Transición a lo Monolingüe}{Juego de adivinar la palabra a través de su definición en inglés simple (\textit{Learner's dictionary}).}
\sesion{24}{Expectativas Mes 3}{Introducción al formato Podcast y preparación psicológica para el ``Club de Escucha''.}
\end{cajasemana}

% ═══════════════════════════════════════════════════════════════════════════════
% MES 3
% ═══════════════════════════════════════════════════════════════════════════════

\newpage

\begin{cajames}[\textcolor{white}{Mes 3: Consolidación Auditiva y Clubes de Debate (A2 $\rightarrow$ B1)}]{mes3}
\textbf{Dinámica:} ``Club de escucha''. El profesor funciona como compañero de conversación, utilizando \textit{recasting} (reformulando en lugar de corregir destructivamente).
\end{cajames}

\vspace{0.3cm}

\begin{cajasemana}[\textcolor{white}{Semana 9: Inmersión Auditiva Pura}]{mes3}
\sesion{25}{Club de Escucha I}{Reproducción grupal de un Podcast (Ej. \textit{Easy Stories in English}). Segmentación de frases.}
\sesion{26}{Debate del Podcast}{Preguntas de comprensión simples. Los alumnos responden con frases cortas. El profesor reformula.}
\sesion{27}{Minería Avanzada (Modismos)}{Revisión en grupo de \textit{collocations} e idioms detectados en el podcast.}
\end{cajasemana}

\begin{cajasemana}[\textcolor{white}{Semana 10: Producción Temprana Guiada}]{mes3}
\sesion{28}{Club de Escucha II}{Tema no ficción/cultural. Entrenar la escucha sin apoyo visual.}
\sesion{29}{Debate Abierto}{Los alumnos comparten opiniones sobre el episodio. Mayor énfasis en conectar ideas (\textit{because, but, so}).}
\sesion{30}{Peer-Review de Minería}{Intercambio de estructuras complejas (B1) minadas asincrónicamente.}
\end{cajasemana}

\begin{cajasemana}[\textcolor{white}{Semana 11: Acentos Diversos y Actualidad}]{mes3}
\sesion{31}{Club de Escucha III}{Noticias adaptadas (Ej. VOA Learning English).}
\sesion{32}{Debate de Actualidad}{Discusión simplificada sobre la noticia escuchada.}
\sesion{33}{Exposición a Acentos}{Sesión dedicada a escuchar clips de acentos globales (Irlandés, Australiano) para ganar agilidad.}
\end{cajasemana}

\begin{cajasemana}[\textcolor{white}{Semana 12: Independencia (B1)}]{mes3}
\sesion{34}{Club de Debate Final}{Tema elegido colaborativamente por los estudiantes. Simulación de charlas de café.}
\sesion{35}{Role-play de Situaciones Reales}{Viajes, compras, entrevistas sencillas. Comprobación pragmática del nivel B1.}
\sesion{36}{Cierre y Graduación}{Retrospectiva del progreso de los 3 meses. Guía final del profesor para continuar la inmersión hacia B2+ de manera 100\% independiente.}
\end{cajasemana}

\vfill
\begin{center}
\rule{0.5\textwidth}{0.5pt}\\[0.3cm]
{\small\color{grissuave} Cronograma Sincrónico · 36 Sesiones · 12 Semanas · Curso de Inmersión en Inglés}
\end{center}

\end{document}
